\documentclass[11pt,a4paper]{article}

\usepackage{fontspec}
\usepackage[english,russian]{babel}
\usepackage{amsmath}
\usepackage{unicode-math}
\setmainfont{Cambria}
\setsansfont{Calibri}
\setmonofont{Consolas}[Scale=0.88]
\setmathfont{Cambria Math}

\usepackage[margin=2.0cm]{geometry}
\usepackage{graphicx}
\usepackage{booktabs}
\usepackage{tabularx}
\usepackage{multirow}
\usepackage{float}
\usepackage{xcolor}
\usepackage{enumitem}
\usepackage[most]{tcolorbox}
\usepackage[font=small,labelfont=bf]{caption}
\usepackage[colorlinks=true,linkcolor=accent,urlcolor=accent,citecolor=accent]{hyperref}

\definecolor{accent}{HTML}{2F6DB5}
\definecolor{good}{HTML}{2E7D4F}
\definecolor{bad}{HTML}{B23A48}
\newcommand{\pos}[1]{\textcolor{good}{#1}}
\newcommand{\negv}[1]{\textcolor{bad}{#1}}
\newtcolorbox{keybox}{colback=accent!5,colframe=accent!60,boxrule=0.6pt,arc=2pt,left=8pt,right=8pt,top=6pt,bottom=6pt}
\setlist{itemsep=2pt,topsep=3pt}
\setlength{\parskip}{4pt}
\setlength{\parindent}{0pt}
\renewcommand{\arraystretch}{1.12}
\emergencystretch=3em
\sloppy

\newcommand{\fc}{\mathrm{fc}}
\DeclareMathOperator*{\argmax}{arg\,max}

\title{\textbf{Устаревание драфта EAGLE после дообучения target}\\[4pt]
\large Карта, предсказание по весам, механизм и дешёвая починка без генерации}
\author{}
\date{15 октября 2026}

\begin{document}
\maketitle
\vspace{-1.0cm}

\begin{keybox}
\textbf{Коротко}
\begin{enumerate}[leftmargin=*]
\item \textbf{Проблема.} Официальный драфт EAGLE-3 для Llama-3.1-8B-Instruct на 9 публичных дообучениях этой модели
теряет \textbf{0--10\%} τ после LoRA, \textbf{10--26\%} после полного дообучения по домену и \textbf{40--60\%} после
дообучения под рассуждения: ускорение падает с 2,34× до 1,5×.
\item \textbf{Предсказание без генерации.} Потерю видно по одним весам: среднее относительное изменение матриц,
пишущих в остаточный поток, коррелирует с ней с \textbf{r = 0,98} (ρ = 0,87, 9 моделей). По сдвигу скрытых состояний
--- r = 0,93.
\item \textbf{Механизм.} «Драфт сам подстраивается, потому что читает признаки target» --- неверно в чистом виде.
Новые признаки \textbf{одновременно} передают драфту новое поведение (+6…+23 п.п. там, где target изменил ответ)
и \textbf{искажают} его работу там, где target ответ не менял (до −19 п.п. у тяжёлых изменений). Искажение идёт на
75--80\% через последний читаемый слой и сосредоточено в 1--4 направлениях.
\item \textbf{Метод.} Поправка ранга 2 на слой, встроенная в $W_{\fc}$ драфта и подобранная за ~1 минуту по прямым
проходам двух моделей (без генерации новым target и без обучения драфта): \textbf{+8,6\%} и \textbf{+9,6\%} τ дерева
на двух моделях с рассуждениями (95\% интервалы выше нуля), +0,2\% на контроле, но \negv{−9,7\%} на abliterated.
\end{enumerate}
\end{keybox}

\section{Постановка}

\textbf{Драфт EAGLE-3.} Target с параметрами $\theta$ на позиции $t$ даёт скрытые состояния $h^{(\ell)}_t(\theta)$.
Драфт читает три слоя $\ell\in\{2,16,29\}$ (входы слоёв) и сжимает их линейным слоем
\[
u_t = W_{\fc}\,\big[h^{(2)}_t;\,h^{(16)}_t;\,h^{(29)}_t\big],\qquad W_{\fc}\in\mathbb R^{4096\times 12288}.
\]
Далее один слой декодера (u поступает и в нормированное внимание, и напрямую в остаточный поток драфта) и голова на
32\,000 частых токенов. Драфт предлагает дерево из 60 токенов глубины 7, target проверяет его за один проход.
$\tau$ --- среднее число принятых токенов за проход target. Генерация жадная, поэтому без потерь качества.

\textbf{Устаревание.} Target дообучают: $\theta_0\to\theta_1$, драфт $\phi_0$ остаётся прежним. Потеря на области $d$:
\[
L_d(\theta_1)=1-\frac{\tau_d(\phi_0;\theta_1)}{\tau_d(\phi_0;\theta_0)} .
\]
\textbf{Вопросы.} (1) насколько велика $L$ на реальных дообучениях; (2) можно ли предсказать её без генерации;
(3) почему драфт не подстраивается; (4) можно ли починить драфт без генерации новым target и без обучения.

\textbf{Что пробовали раньше (7--13 октября).} На LoRA (Qwen3-8B, Llama-3.1-8B) поправка драфта под конкретное
изменение оказалась малой (+0,03…+0,10 τ) по сравнению с общей докалибровкой; Procrustes, ridge и малые базисы её не
воспроизводили ($\le$5\% действия). Вывод: LoRA слишком слабо меняет target, нужны тяжёлые изменения --- отсюда карта.

\section{Карта устаревания}

\textbf{Протокол.} Код авторов EAGLE (коммит \texttt{cb7e084}), официальный драфт \texttt{yuhuili/EAGLE3-LLaMA3.1-Instruct-8B},
дерево 60/7/10, жадно, до 512 новых токенов; 96 запросов: по 16 из общих вопросов, математики, кода, медицины,
вызовов инструментов (ToolACE) и русского (m-ArenaHard). Colab A100 80GB. Ускорение --- на 12 запросах.

\begin{table}[H]
\centering\footnotesize
\begin{tabular}{@{}llcccccc c@{}}
\toprule
Target & тип изменения & общее & матем. & код & медиц. & инстр. & рус. & ускор. \\
\midrule
Llama-3.1-8B-Instruct & --- & 5,71 & 6,13 & 6,44 & 5,15 & 3,95 & 1,67 & 2,34 \\
\midrule
LIMO & LoRA r8 & 5,31 & 5,54 & 6,42 & 4,98 & 3,82 & 1,64 & 1,93 \\
ToolACE-8B & полное (≈ LoRA) & 4,96 & 5,23 & 5,87 & 5,09 & 3,65 & 1,59 & 2,09 \\
SuperNova-Lite & дистилляция & 5,03 & 5,12 & 5,91 & 4,85 & 3,14 & 1,47 & 1,87 \\
Vikhr (рус.) & полное & 4,63 & 5,17 & 5,84 & 4,61 & 3,39 & 1,53 & 2,02 \\
UltraMedical & полное & 4,26 & 5,02 & 5,42 & 4,45 & 4,05 & 1,45 & 1,87 \\
HuatuoGPT-o1 & полное, рассуждения & 4,45 & 4,51 & 4,84 & 4,14 & 4,28 & 3,41$^*$ & 2,87 \\
abliterated & тяжёлое & 3,38 & 3,41 & 4,36 & 4,02 & 2,87 & 2,27 & 2,24 \\
R1-Distill & дистилляция рассуждений & 2,82 & 3,12 & 3,00 & 3,20 & 2,83 & 1,35 & 1,74 \\
Nemotron-Nano & рассуждения, SFT+RL & 2,49 & 2,60 & 2,55 & 2,96 & 2,30 & 1,57 & 1,51 \\
\midrule
R1-Distill, \emph{свой} драфт & --- & 6,16 & 7,34 & 6,13 & 5,13 & 2,09 & 1,47 & 2,37 \\
Instruct, драфт R1 & --- & 2,05 & 2,58 & 2,35 & 2,19 & 1,77 & 1,34 & 1,30 \\
\bottomrule
\end{tabular}
\caption{τ дерева. $^*$HuatuoGPT-o1 отвечает на русские вопросы по-английски (7\% кириллицы). На русском у исходной
модели 56\% токенов ответа вне словаря драфта и 93\% обрывов цепочки происходят на них.}
\end{table}

\begin{figure}[H]
\centering\includegraphics[width=0.78\textwidth]{fig/fig_pitch_1_ladder.png}
\caption{Средняя потеря τ по 5 английским областям (столбцы) и по областям (точки).}
\end{figure}

\textbf{Выводы.} LoRA: 0--10\%; полное дообучение по домену: 1--26\% во \emph{всех} областях, не только в «своей»;
тяжёлые изменения и рассуждения: 22--60\%. Собственный драфт возвращает всё (2×2 с R1-Distill), значит потеря ---
несоответствие драфта, а не свойство модели. Обучающие данные драфта тоже важны: на вызовах инструментов собственный
драфт R1-Distill хуже драфта Instruct (2,09 против 2,83).

\section{Предсказание устаревания без генерации}

Для каждого изменения на общих текстах (ответы исходной модели на 48 запросов) считаем
\begin{gather*}
\delta=\mathbb E_t\,\frac{\|h_t(\theta_1)-h_t(\theta_0)\|}{\|h_t(\theta_0)\|}\quad\text{(прямые проходы)},\\
\rho_W=\operatorname{mean}_{W\in\{o,\;down\}}\frac{\|W_1-W_0\|_F}{\|W_0\|_F}\quad\text{(только веса)} .
\end{gather*}

\begin{table}[H]
\centering\footnotesize
\begin{tabular}{@{}lccc@{}}
\toprule
Предсказатель потери τ (9 изменений) & Пирсон [95\%] & Спирмен & ошибка LOO \\
\midrule
$\rho_W$, веса o\_proj/down\_proj (\textbf{только веса}) & \textbf{0,98} & 0,87 & --- \\
сдвиг состояний $\delta$, три слоя & 0,93 [0,55; 0,98] & 0,83 & 7,7 п.п. \\
сдвиг состояний, слой 29 & 0,94 [0,65; 0,99] & 0,90 & 7,1 п.п. \\
KL драфта (новые против старых признаков) & 0,93 [0,78; 0,99] & 0,90 & 6,6 п.п. \\
изменение поведения target (доля argmax) & 0,86 & 0,78 & 9,4 п.п. \\
константа & --- & --- & 13,8 п.п. \\
\bottomrule
\end{tabular}
\end{table}

\begin{figure}[H]
\centering\includegraphics[width=0.95\textwidth]{fig/fig_pitch_2_predictor.png}
\end{figure}

\textbf{Оговорки.} Корреляцию держат три тяжёлых изменения; без двух моделей с рассуждениями r = 0,75 (ρ = 0,64).
Умеренные дообучения точно не ранжируются: HuatuoGPT-o1 меняет веса меньше чем на 1\%, но теряет 17\%. Это
фильтр «сломается ли драфт сильно», а не точный прогноз. Сдвиг состояний предсказывает лучше, чем изменение
поведения самого target.

\textbf{Как меняются веса.} LIMO и ToolACE --- низкоранговые изменения (32 направления несут 89\% и 99,5\%
энергии $\Delta W$); остальные плотные (5--58\%). Направления, в которых веса пишут в остаточный поток (по 8 главных
левых сингулярных векторов $\Delta W_{o},\Delta W_{down}$ предыдущих слоёв), содержат \textbf{21--66\%} среднего
сдвига состояний в слое 29 против 7--13\% у случайного подпространства той же размерности: веса подсказывают, куда
сдвинутся признаки. R1-Distill и Nemotron-Nano (разные команды, разные стартовые точки) сдвигают состояния в общем
направлении: косинус средних сдвигов 0,5--0,86, устойчиво к исключению «массивных» измерений.

\section{Механизм: перенос против искажения}

На общем тексте делим позиции по тому, изменил ли target свой следующий токен: $y^1_t=y^0_t$ (множество $S$) или
$y^1_t\ne y^0_t$ (множество $C$). Совпадение драфта на первом шаге, когда он читает признаки $f$:
\[
a_X(f)=\Pr_{t\in X}\big[\argmax q_{\phi_0}(\cdot\mid f_{\le t})=y^1_{t}\big],\qquad
\underbrace{T=a_C(f^1)-a_C(f^0)}_{\text{перенос}},\qquad
\underbrace{D=a_S(f^0)-a_S(f^1)}_{\text{искажение}} .
\]

\begin{table}[H]
\centering\small
\begin{tabular}{@{}lccccccccc@{}}
\toprule
 & LIMO & ToolACE & SuperN. & Vikhr & UltraM. & Huatuo & ablit. & R1-D. & Nemotron \\
\midrule
доля $C$, \% & 5,1 & 4,5 & 3,3 & 7,3 & 8,4 & 7,2 & 7,2 & 16,4 & 14,5 \\
перенос $T$, п.п. & \pos{+14,5} & \pos{+18,2} & \pos{+12,2} & \pos{+18,0} & \pos{+20,4} & \pos{+22,7} & \pos{+17,0} & \pos{+8,5} & \pos{+5,6} \\
искажение $D$, п.п. & 0,3 & 0,2 & 0,2 & 0,8 & 1,2 & 1,4 & \negv{3,4} & \negv{10,0} & \negv{19,3} \\
\bottomrule
\end{tabular}
\end{table}

\begin{figure}[H]
\centering\includegraphics[width=0.78\textwidth]{fig/fig_pitch_3_mechanism.png}
\end{figure}

\begin{itemize}[leftmargin=*]
\item У умеренных изменений потеря --- это сами изменения поведения ($C$: 3--8\% токенов, драфт угадывает их лишь в
~45\% случаев); признаки драфт почти не портят.
\item У тяжёлых изменений добавляется искажение на $S$ --- 85--95\% токенов, где ответ target прежний. Оно бьёт по
частым служебным токенам (\texttt{the}, запятая, точка), а не по содержанию, и затрагивает все токены.
\item \textbf{Не масштаб, а направление.} У рассуждающих моделей $\cos(h^0,h^1)=0{,}68$--$0{,}81$, чистое
масштабирование объясняет 2--15\% сдвига. Через $W_{\fc}$ драфт видит сдвиг на 73--80\% из слоя 29, на ~25\% из
слоя 16 и почти не видит из слоя 2.
\item Отсюда старая формулировка «бесплатного переноса нет»: в сумме по токенам $T$ и $D$ компенсируют друг друга.
\end{itemize}

\section{Метод: избирательное восстановление входа драфта}

Цель --- убрать искажение $D$, сохранив перенос $T$. Восстановление строится как функция только новых признаков и
подгоняется по прямым проходам $\theta_0$ и $\theta_1$ на 24 текстах исходной модели (половина общих текстов; λ ---
на отложенной четверти).

\textbf{Поправка (зафиксирована до проверки τ).} Для каждого читаемого слоя $j$: средний сдвиг $\mu_j$ и главное
направление $v_j$ центрированного сдвига, $P_j=\mathrm{orth}[\mu_j;v_j]\in\mathbb R^{2\times4096}$. Ridge-регрессия
предсказывает, насколько вернуть координаты вдоль $P_j$:
\[
\min_{R_j}\big\|(X^0-X^1)P_j^\top-[X^1,\mathbf 1]R_j\big\|_F^2+\lambda\|R_j\|_F^2,\qquad
x'=x+[x,1]R_jP_j=x\,(I+A_j)+c_j ,
\]
где $A_j=R_j^{(1:H)}P_j$ (ранг 2), $c_j=R_j^{(H+1)}P_j$. Встраиваем в драфт:
\[
W_{\fc}'=W_{\fc}\,\mathrm{blockdiag}\big(I+A_2^\top,\,I+A_{16}^\top,\,I+A_{29}^\top\big),\qquad
b_{\fc}=W_{\fc}\,[c_2;c_{16};c_{29}] .
\]
Стоимость: прямые проходы двух моделей на ~10 тыс. токенов и три решения линейных систем, около минуты на A100.
Новый target ничего не генерирует, драфт не обучается.

\textbf{Выбор варианта (первый шаг драфта, тестовая половина).} 30 вариантов: среднее, поэлементное аффинное,
$k\in\{1,4,16,64\}$ направлений сдвига, разные наборы слоёв. Лучшие возвращают 47--57\% разрыва
$a(f^0)-a(f^1)$; лучшим оказывается 1--4 направления, а не 64 --- искажение почти одномерно. Перенос при этом
сохраняется: у Nemotron на $C$ совпадение 0,281 → 0,320 (со старыми признаками 0,204).

\begin{table}[H]
\centering\small
\begin{tabular}{@{}lcccc@{}}
\toprule
Изменение τ дерева & общее начало & 95\% интервал & запросов & по областям \\
с поправкой & ответов & & лучше / хуже & \\
\midrule
Nemotron-Nano & \pos{\textbf{+8,6\%}} & [+7,4; +10,0] & 94\% / 3\% & +6,5…+11,2\% \\
R1-Distill & \pos{\textbf{+9,6\%}} & [+8,4; +10,8] & 90\% / 10\% & +3,9…+10,7\% \\
ToolACE (контроль) & +0,2\% & [−0,5; +1,0] & 24\% / 33\% & −2,3…+4,8\% \\
abliterated & \negv{−9,7\%} & [−11,0; −8,3] & 3\% / 89\% & −12,2…−7,8\% \\
\bottomrule
\end{tabular}
\caption{Протокол карты, ответы самого изменённого target. Тексты с поправкой и без не всегда побуквенно одинаковы
(bf16-проверка дерева иначе разрешает почти равные варианты), поэтому сравнение --- на общем начале ответов до первого
расхождения (43--80\% длины). Русский: без изменений.}
\end{table}

\begin{figure}[H]
\centering\includegraphics[width=0.62\textwidth]{fig/fig_pitch_4_fix.png}
\end{figure}

Выигрыш идёт за счёт меньшей доли пустых циклов (Nemotron 23\% → 20\%) и большего числа длинных цепочек, равномерно
по ответу. При неизменной стоимости цикла ускорение растёт 1,51× → ~1,65× (Nemotron) и 1,74× → ~1,86× (R1-Distill).

\textbf{Почему abliterated хуже.} Там искажение в основном \emph{масштабное}: $\|u^1\|/\|u^0\|=0{,}79$, поворот
слабый ($\cos(u^0,u^1)=0{,}94$), 45\% сдвига состояний --- чистое масштабирование. Выход $\fc$ идёт в остаточный поток
драфта без нормировки, так что масштаб важен, а поправка рассчитана на поворот. Прокси на первом шаге обещал здесь лишь
+0,011 и не увидел вреда. \textbf{Гипотеза правила:} включать поправку при медиане $\cos(u^0,u^1)<0{,}9$ ---
отделяет помогающие случаи (0,75; 0,82) от остальных (0,94; 0,99). Подобрано задним числом по 4 точкам.

\section{Новизна и связанные работы}

\begin{itemize}[leftmargin=*]
\item \textbf{EDA} (2026): адаптация драфта к дообученному target через генерацию и обучение; у нас --- без генерации и
обучения. \textbf{AngelSpec}, \textbf{Draft-OPD}: драфт режима без рассуждений хуже на режиме с рассуждениями (смена
режима, не весов). \textbf{SPEED-Bench}, \textbf{AdaSpec}, \textbf{SpecVocab}: урезанный словарь вредит неанглийским
языкам --- наши цифры по русскому это подтверждают.
\item Не найдено в литературе: предсказание устаревания драфта по весам изменения; разложение «перенос против
искажения» с прямым замером; починка входа драфта из прямых проходов без генерации.
\end{itemize}

\section{Ограничения и следующие шаги}

\begin{itemize}[leftmargin=*]
\item Одна база (Llama-3.1-8B) и один драфт; 16 запросов на область; метод проверен на 4 изменениях, правило включения
--- задним числом. Ускорение по времени оценено через τ, а не измерено.
\item \textbf{Следующие шаги:} (1) заранее объявленная проверка на 2--3 новых моделях с рассуждениями с правилом
$\cos<0{,}9$ и замером скорости (~1,5 ч A100); (2) вторая семья моделей (Qwen3-8B); (3) поправка, учитывающая масштаб
(для abliterated); (4) веса как источник направлений поправки --- шаг к версии «только по весам».
\end{itemize}

\end{document}
